\chapter*{Prakata}
\addcontentsline{toc}{chapter}{Prakata}

Tujuan buku ini adalah menyajikan pembahasan yang
ketat,
konservatif,
elementer,
dan minimalis.
Pada saat yang sama, isinya mendekati batas maksimum yang dapat dipelajari mahasiswa dalam satu semester.

Buku ini disusun berdasarkan
perkuliahan yang diberikan oleh penulis
di Pennsylvania State University.
Perkuliahan tersebut ditujukan kepada mahasiswa sarjana tahun kedua dan tahun terakhir yang telah mengenal bilangan real dan kekontinuan.
Bekal ini memungkinkan materi dibahas lebih cepat dan lebih ketat daripada di sekolah menengah.
Sekitar sepertiga materi ini dahulu diajarkan di sekolah menengah, tetapi kini tidak lagi.

\section{Prasyarat}

Mahasiswa sebaiknya telah mengenal
topik-topik berikut:
\begin{itemize}
\item Teori himpunan elementer:
$\in$,
$\cup$, 
$\cap$,
$\backslash$,
$\subset$,~$\times$.
\item Bilangan real: interval, ketaksamaan, identitas aljabar.
\item Limit, fungsi kontinu, dan teorema nilai antara.
\item Fungsi standar:
nilai mutlak,
logaritma natural,
fungsi eksponensial.
Sesekali digunakan pula fungsi trigonometri.
\item Bab~\ref{chap:trans} menggunakan aljabar vektor dasar.
\item Untuk membaca Bab~\ref{chap:sphere}, sebaiknya pembaca telah mengenal \textit{hasil kali skalar}, yang juga disebut \textit{hasil kali titik}.
\item Untuk membaca Bab~\ref{chap:complex}, sebaiknya pembaca telah mengenal bilangan kompleks.
\end{itemize} 

\section{Gambaran umum}

Kita menggunakan apa yang disebut \textit{pendekatan metrik}, yang diperkenalkan oleh Birkhoff.
Artinya, kita mendefinisikan bidang Euklides sebagai \textit{ruang metrik} yang memenuhi suatu daftar sifat (\textit{aksioma}).
Dengan cara ini, kita meminimalkan bagian-bagian teknis yang menjemukan
dan tidak dapat dihindari dalam pendekatan Hilbert yang lebih klasik.
Pada saat yang sama, mahasiswa berkesempatan mempelajari geometri ruang metrik.

Berikut adalah graf ketergantungan antarbab.

\begin{figure}[!ht]
\centering

\begin{tikzpicture}[->,>=stealth',shorten >=1pt,auto,scale=1.4,
  thick,main node/.style={circle,draw,font=\sffamily\bfseries,minimum size=8mm}]

  \node[main node] (1) at (-1.5,10/6) {\ref{chap:metr}};
  \node[main node] (2) at (-.5,10/6){\ref{chap:axioms}};
  \node[main node] (3) at (.5,10/6) {\ref{chap:half-planes}};
  \node[main node] (4) at (1.5,10/6) {\ref{chap:cong}};
  \node[main node] (5) at (2.5,10/6) {\ref{chap:perp}};
  \node[main node] (6) at (3.5,10/6) {\ref{chap:parallel}};
  \node[main node] (7) at (4.5,10/6) {\ref{chap:angle-sum}};
  \node[main node] (8) at (5.5,10/6) {\ref{chap:triangle}};
  \node[main node] (9) at (5,5/6){\ref{chap:inscribed-angle}};
  \node[main node] (10) at (4.5,0) {\ref{chap:inversion}};
  \node[main node] (11) at (3,5/6) {\ref{chap:non-euclid}};
  \node[main node] (12) at (3.6,0){\ref{chap:poincare}};
  \node[main node] (13) at (3,-5/6) {\ref{chap:h-plane}};
  \node[main node] (14) at (5,-5/6) {\ref{chap:trans}};
  \node[main node] (15) at (4.5,-10/6) {\ref{chap:proj}};
  \node[main node] (16) at (4,-5/6) {\ref{chap:sphere}};
  \node[main node] (17) at (3.5,-10/6) {\ref{chap:klein}};
  \node[main node] (18) at (5.5,0) {\ref{chap:complex}};
  \node[main node] (19) at (4,15/6) {\ref{chap:car}};
  \node[main node] (20) at (5,15/6) {\ref{chap:area}};

  \path[every node/.style={font=\sffamily\small}]
   (1) edge node{}(2)
   (2) edge node{}(3)
   (3) edge node{}(4)
   (4) edge node{}(5)
   (5) edge node{}(6)
   (6) edge node{}(7)
   (7) edge node{}(9)
   (7) edge node{}(8)
   (7) edge node{}(20)
   (10) edge node{}(14)
   (14) edge node{}(15)
   (7) edge node{}(19)
   (9) edge node{}(10)
   (10) edge node{}(12)
   (10) edge node{}(16)
   (10) edge node{}(18)
   (5) edge node{}(11)
   (11) edge node{}(12)
   (12) edge node{}(13)
   (15) edge node{}(17)
   (16) edge[dashed] node{}(17)
   (13) edge node{}(17);
\end{tikzpicture}
\end{figure}

Dalam (\ref{chap:metr}), kita memperkenalkan konsep-konsep yang diperlukan untuk merumuskan aksioma;
bahasannya mencakup ruang metrik,
garis,
ukuran sudut,
pemetaan kontinu,
dan segitiga kongruen.

Selanjutnya, kita memasuki geometri Euklides:
(\ref{chap:axioms}) Aksioma dan akibat langsung;
(\ref{chap:half-planes}) Setengah-bidang dan kekontinuan;
(\ref{chap:cong}) Segitiga kongruen;
(\ref{chap:perp}) Lingkaran, transformasi isometrik, dan garis tegak lurus;
(\ref{chap:parallel}) Segitiga sebangun dan (\ref{chap:angle-sum}) Garis sejajar
--- inilah dua bab pertama tempat kita menggunakan Aksioma~\ref{def:birkhoff-axioms:4}, yang ekuivalen dengan postulat kesejajaran Euklides.
Dalam (\ref{chap:triangle}), kita menyajikan teorema-teorema geometri segitiga yang paling klasik;
bab ini terutama disertakan sebagai ilustrasi.

Dalam dua bab berikutnya, kita membahas geometri lingkaran pada bidang Euklides:
(\ref{chap:inscribed-angle}) Sudut keliling; (\ref{chap:inversion}) Inversi.
Inversi akan digunakan untuk membangun model bidang hiperbolik.

Selanjutnya, kita membahas geometri non-Euklides:
(\ref{chap:non-euclid})
Geometri netral --- geometri tanpa postulat kesejajaran;
(\ref{chap:poincare})
Model cakram konformal ---
suatu konstruksi bidang hiperbolik,
yakni contoh bidang netral yang bukan bidang Euklides.
Dalam (\ref{chap:h-plane}), kita membahas geometri bidang hiperbolik yang telah dibangun --- inilah puncak pembahasan buku ini.

Dalam bab-bab sisanya, kita membahas topik tambahan:
(\ref{chap:trans}) Geometri afin;
(\ref{chap:proj}) Geometri proyektif;
(\ref{chap:sphere}) Geometri sferis;
(\ref{chap:klein}) Model proyektif bidang hiperbolik;
(\ref{chap:complex}) Koordinat kompleks;
(\ref{chap:car}) Konstruksi geometri;
(\ref{chap:area}) Luas.
Bukti-bukti dalam bab-bab ini tidak sepenuhnya ketat.

Kami menganjurkan pembaca menggunakan tugas visual pada situs web penulis.

\textbf{Latihan kelas} biasanya mengilustrasikan konsep yang baru diperkenalkan dan dimaksudkan untuk dibahas di kelas.
Latihan yang digunakan kembali pada bagian selanjutnya ditandai dengan tanda seru: \textbf{Latihan!}
Suatu pernyataan ditandai dengan ``$\a$''\label{a-mark} (misalnya, ``\textbf{Teorema\abs}'') jika buktinya tidak menggunakan Aksioma~\ref{def:birkhoff-axioms:4}.
Penandaan ini akan menjadi penting mulai Bab~\ref{chap:non-euclid}.


\section{Penafian}
 
Sebagian besar bukti dalam buku ini telah muncul dalam Elements karya Euklides.
Saya tidak berusaha mencari rujukan aslinya; tugas itu mustahil dilakukan.

\section{Sumber yang disarankan}

\begin{itemize}
\item {}\emph{Byrne's Euclid} \cite{byrne} --- versi berwarna dari enam buku pertama Elements karya Euklides, disunting oleh Oliver Byrne.

\item {}\emph{Kiselyov's geometry} \cite{kiselev} ---
buku teks klasik untuk pelajar sekolah karya Andrey Kiselyov; buku ini dapat membantu jika Anda kesulitan mengikuti pembahasan di sini.

\item {}\emph{Lessons in Geometry} karya Jacques Hadamard \cite{hadamard} --- ensiklopedia geometri elementer yang semula ditulis untuk guru sekolah.

%\item Moise's book, \cite{moise} --- should be good for further study.

%\item Greenberg's book \cite{greenberg}  --- a historical tour in the axiomatic systems of various geometries.

\item {}\emph{Problems in geometry} karya Victor Prasolov \cite{prasolov} sempurna untuk menguasai keterampilan pemecahan masalah.

\item {}\emph{Geometry in figures} karya Arseniy Akopyan \cite{akopyan} --- ensiklopedia geometri Euklides yang nyaris tanpa kata-kata.

\item {}\emph{Euclidea} \cite{euclidea} --- cara yang menyenangkan sekaligus menantang untuk mempelajari konstruksi geometri.
 
\item {}\emph{Geometry} karya Igor Sharygin \cite{sharygin} --- buku teks geometri terbaik bagi pelajar sekolah.
Saya merekomendasikannya kepada siapa pun yang dapat membaca bahasa Rusia.

\item {}\emph{problems.ru} \cite{zadachi} --- kumpulan soal daring dalam bahasa Rusia.

\end{itemize}

\section{Ucapan terima kasih}

{\sloppy

Saya berterima kasih kepada
Stephanie Alexander,
Thomas Barthelme,
Berk Ceylan,
Matthew Chao, 
Quinn Culver,
Svetlana Katok, 
Nina Lebedeva,
Alexander Lytchak,
Alexei Novikov,
dan Lukeria Petrunina
atas saran-saran yang berguna dan koreksi terhadap kesalahan cetak.

Karya ini didukung sebagian oleh
hibah NSF
DMS-0103957,
DMS-0406482,
DMS-0905138,
DMS-1309340,
DMS-2005279,
serta hibah Simons Foundation
245094, 584781.

}
